\section{Results}
\label{sec:results}

Table~\ref{tab:main-results} summarizes the principal outcomes.  The values are
exact for the fixed three-case corpus and deterministic test paths.  They are
not participant counts or dialogue-quality scores.

\begin{table*}[t]
  \caption{Main measured results and their interpretation boundary.}
  \label{tab:main-results}
  \small
  \begin{tabularx}{\textwidth}{@{}p{31mm}Y Y@{}}
    \toprule
    Check & Observed result & Supported interpretation \\
    \midrule
    Complete declarations
      & 93/93 scene assets and 186/186 asset-targeting
        step--capability-use pairs complete
      & The three regenerated models contain provider-coherent paths through
        binding, action, assertion, and validation coverage. \\
    Migration non-regression
      & 93/93 owner decisions and 459/459 route classifications agree
      & V2 preserves the direct artifacts under the fixed shared resolver; this
        is not an independent routing oracle. \\
    Runtime admission
      & 298/298 local/direct routes admitted; 161/161
        transitive/unreachable routes rejected before the stub
      & The implemented one-hop policy is enforced and preserves target and
        provider evidence without retained mutation on rejection. \\
    Common injected faults
      & Both adapters detected and localized 9/9; 0/3 false positives
      & Parity for the three shared fault types. \\
    Contract-only faults
      & Strict provider contract detected and localized 9/9; 0/3 false
        positives; canonical model validator detected 3/9
      & The executable profile adds typed cross-layer invariants absent from
        the direct artifacts. \\
    Unity execution
      & 3/3 batch entry points passed; one multi-scenario smoke loaded 31
        contexts and 73 actions and bound 30/30 scene objects exactly
      & The evaluated desktop runtime components execute the generated
        contract; no headset or user outcome is established. \\
    \bottomrule
  \end{tabularx}
\end{table*}

\subsection{RQ1: complete migration with routing non-regression}

All 93 scene assets occurred in at least one declared interaction chain, and
all 93 satisfied the asset-completeness definition.  The chain audit examined
186 asset-targeting \texttt{ScenarioStep}/\texttt{CapabilityUse} pairs:
54 in the career-fair case, 72 in the classroom, and 60 in the brand room.
Every pair resolved one capability and provider, agreed with its scenario
performer and target, reached an executable binding and action, resolved all
resulting assertions, and was covered by an appropriate validation case.
No asset-chain error was reported.  The two pairs per asset encode answer and
handoff paths; they are not two executed conversations.

Natural-corpus ownership was unique for 93/93 assets.  Every decision occurred
at the explicit asset tier; no natural object exercised group or zone fallback,
and none was ambiguous or unassigned.  In the separate derived priority suite,
both adapters passed all nine decision cases: three group fallbacks, three zone
fallbacks, and three ambiguity rejections.  Across both adapters this is 18/18
expected outcomes, kept separate from natural-corpus frequencies.

The 459 object-by-start-agent routes comprised 93 local-owner cases, 205 direct
handoffs, 143 transitive-only graph paths, and 18 unreachable pairs.
Local-or-direct one-hop coverage was therefore 298/459 (64.9\%), while offline
graph reachability was 441/459 (96.1\%).  The latter does not imply that the
online runtime traverses several handoffs within one turn.

The direct-wiring and fresh-V2 adapters selected the same owner for all 93
objects and the same route class for all 459 pairs.  They also remained aligned
after all nine common mutations.  Because both normalized views are classified
by one shared policy function, this result establishes migration
non-regression and serialization parity, not that V2 discovered a more correct
routing policy.

\subsection{RQ2: fail-closed enforcement}

The executable runtime sweep passed all 459 probes.  It admitted 298/298 local
or direct routes and made exactly 298 deterministic-stub calls.  Every admitted
response preserved the expected source object, target entity, and routed
provider and returned non-empty capability-use, capability, binding, and action
identifiers that agreed across the checked response, grounding, routing, and
runtime-action fields.

The runtime rejected 161/161 transitive or unreachable routes before the stub
and without retained session mutation.  This set contains the 143
transitive-only and 18 unreachable pairs.  The result demonstrates the current
one-hop admission boundary, not task failure in principle: a system with a
different declared policy could permit multi-hop execution, but this
implementation does not.

The targeted spatial suite accepted its valid deictic request and rejected all
16 invalid variants without retained mutation.  Fifteen variants were
request-decidable and stopped before the stub.  The remaining variant used a
structured response that proposed an undeclared handoff; one stub call was
necessary to reveal the destination, after which provisional history was
rolled back.  Rejections included stale model hashes, unknown or contradictory
identifiers, ambiguous selections, invalid agents, excessive payloads, and
non-finite or out-of-range spatial values.

All three Unity Editor batch entry points returned success.  Together they
exercised scenario progression, dispatch, five assertion types, valid and
deliberately broken evidence, selection-state blocking, and ambiguous
registries.  The multi-scenario Steinpilz smoke loaded 31 contexts and 73
actions and established unique exact bindings for 30/30 scene objects without
UUID-suffix fallback.  These observations confirm that the client-side bridge
and evidence components execute the materialized contract; they do not measure
frame rate, end-to-end latency, or human understanding.

\subsection{RQ3: added invariant scope and cost}

On the common mutation denominator, both adapters accepted all three unmodified
case copies, for 0/3 false positives per adapter.  Both detected and localized
9/9 changes: one missing owner, one duplicate owner, and one dangling handoff
per scene.  The result is deliberately neutral.  Direct wiring can enforce
these three relationships when supplied with a validator and the same policy.

The representation patches expose a mixed trade-off.  Across the nine common
mutations, the contract representation changed 9 top-level artifacts compared
with 15 for direct wiring, but changed 18 stored references compared with 15.
These counts describe the scripted patches and do not establish that a
developer would work faster or make fewer mistakes.

The contract-only suite provides the positive evidence for additional
expressiveness.  Table~\ref{tab:v2-only} reports the three typed mutation
classes.  The canonical model validator detected the duplicated source-agent
identifier in each case (3/9 overall) but intentionally permits some optional
relations for compatibility.  The executable pipeline agent--provider contract
detected and localized all 9/9 mutations with 0/3 false positives on valid
copies.  The direct artifacts have no typed capability-use provider or
domain-capability-to-runtime-orchestrator relation, so there is no equivalent
baseline outcome to score.

\begin{table*}[t]
  \caption{Contract-only mutations.  ``Not represented'' means that the direct
  artifacts contain no equivalent typed relation; it is not counted as a
  baseline miss.}
  \label{tab:v2-only}
  \small
  \begin{tabularx}{\textwidth}{@{}p{38mm}Y p{26mm}p{30mm}@{}}
    \toprule
    Mutation (once per scene) & Violated invariant & Direct artifacts &
    Executable provider contract \\
    \midrule
    Remove \texttt{CapabilityUse.provider}
      & An executable capability occurrence has exactly one provider agreeing
        with the performing agent and target responsibility.
      & Not represented & 3/3 detected and localized \\
    Duplicate \texttt{Agent.sourceAgentId}
      & Runtime-facing source-agent identity is unique.
      & No typed equivalent & 3/3 detected and localized \\
    Assign a domain capability to the runtime orchestrator
      & Domain behavior is provided by the modeled domain agent, while the
        orchestrator realizes infrastructure actions.
      & Not represented & 3/3 detected and localized \\
    \bottomrule
  \end{tabularx}
\end{table*}

Table~\ref{tab:runtime-cost} reports the current benchmark snapshot.  The
fresh-V2 adapter required 5.89--7.61 times the direct-wiring median for the
measured in-memory workload.  Absolute medians ranged from 6.15 to 10.06\,ms
for V2 and 1.04 to 1.32\,ms for direct wiring.  Outputs remained deterministic
across all repetitions.  These values exclude file parsing, generation,
Unity, networking, and model inference and must not be read as end-user
response latency.

\begin{table}[t]
  \caption{Local structural workload over 40 measured repetitions per adapter
  and case.  Values are median [Q1, Q3] in milliseconds.}
  \label{tab:runtime-cost}
  \small
  \begin{tabular}{@{}lrrr@{}}
    \toprule
    Case & Direct & Fresh V2 & Ratio \\
    \midrule
    Career fair & 1.322 [1.167, 1.516] & 10.057 [7.784, 11.573] & 7.61 \\
    Classroom & 1.246 [0.772, 1.609] & 9.439 [5.681, 13.336] & 7.58 \\
    Brand room & 1.044 [0.992, 1.694] & 6.151 [4.951, 7.309] & 5.89 \\
    \bottomrule
  \end{tabular}
\end{table}

Taken together, RQ3 does not show that an explicit contract is universally
better than direct artifacts.  It shows a more specific trade: shared
ownership and handoff faults can be validated in either representation, while
the contract makes additional provider, identity, capability, binding, and
evidence relations explicit and enforceable at measurable structural cost.

\FloatBarrier
\subsection{RQ4: comprehension of the implemented interaction}
\label{sec:results-userstudy}

The dated user-study cut contains 65 completed sessions.  Target selection was
rated positively (4 or 5) by 45/65 participants (69.2\%; median 4,
IQR 3--5), and 50/65 rated the task-level answer--target fit positively
(76.9\%; median 5, IQR 4--5).  Across the final ratings, 55/65 found the
guidance helpful (84.6\%), 47/65 rated the overall interaction easy (72.3\%),
and the 0--10 overall score had a median of 8 (IQR 7--9).

Forty-nine participants reported observing the intended handoff; eight were
unsure and eight reported no handoff.  Among those 49, 40 rated the transition
understandable (81.6\%; median 5, IQR 4--5) and 45 found the final responding
agent clear (91.8\%; median 5, IQR 5--5).  These values measure perceived
handoff visibility.  The questionnaire export does not contain the
corresponding per-session backend event, so it cannot establish individual
log--report agreement.

Responsibility rationale was less consistently understood.  The general
responsibility item was positive for 42/65 participants (64.6\%; median 4,
IQR 3--4), but only 37/65 selected the intended rule that provider choice
depends on the selected object and question topic (56.9\%; 95\% Wilson CI
[44.8\%, 68.2\%]).  Ten selected the first-contacted agent, eight the nearest
agent, and ten reported that the rule was unclear.  Likewise, only 30/65 rated
the division among several specialized agents as helpful (46.2\%; median 3,
IQR 3--4).  Thus the implemented cues made the handoff event and final agent
more visible than the underlying responsibility rule.

Only 22 participants reported experiencing a state for which they could rate
error explanation; 15/22 rated it positively.  We treat this result as
exploratory because exposure was neither required nor controlled.  Optional
comments primarily identified navigation and input friction, unclear or
overlapping agent roles, and answer-presentation or content-quality issues.
These observations guide interface refinement but do not alter the technical
contract results.

\FloatBarrier
